
%% ============================================================
\appendix
\section{Infrastructure and Pipeline}
\label{sec:appendix-pipeline}
%% ============================================================

%% ============================================================
\subsection{Infrastructure Choices}
%% ============================================================

\subsubsection{Modal --- GPU Compute}

A3 needs single-GPU training jobs. Modal is the simplest platform for this: write a Python function, pick a GPU, run it. We had \$500 in academic credits, per-second billing, and Volumes (persistent networked disks) to save checkpoints between runs.

\paragraph{Why not others}
Anyscale (\$1{,}000) is built around Ray for distributed workloads---more setup than needed for single-GPU jobs. AWS (\$500) has the most overhead (EC2, SSH, drivers, storage). Snowflake (\$400 trial) is a data warehouse, not a training platform.

\paragraph{GPU selection}
Nanochat requires BF16, so A10G is the minimum. A100/H100 are faster and often cheaper in total cost due to shorter wall time.

\begin{table}[H]
\centering\small
\begin{tabular}{@{}llrr@{}}
\toprule
Part & GPU & Est.\ Time & Est.\ Cost \\
\midrule
P2 (4 ablation runs) & A100 & 1--5 hrs & \$3--13 \\
P3 (2 training stages + eval) & A100 & 1--3 hrs & \$3--8 \\
P4 (full nanochat) & H100 & 6--20 hrs & \$24--79 \\
\bottomrule
\end{tabular}
\caption{Estimated GPU time and cost per assignment part.}
\end{table}

\subsubsection{Weights \& Biases --- Experiment Tracking}

Dashboard for training curves, loss values, and run configs. Already built into nanochat---no extra code needed.

\begin{itemize}[nosep]
  \item \textbf{Project}: \texttt{490-autobook-a3} (shared across all parts)
  \item \textbf{Namespacing}: each run tagged \texttt{p2}, \texttt{p3}, or \texttt{p4}
  \item \textbf{Reproducibility}: every run logs which nanochat commit it used
\end{itemize}

\subsubsection{YAML Configs --- Run Specification}

Each experiment is a YAML file whose fields map directly to nanochat's command-line flags (\texttt{--depth}, \texttt{--max-seq-len}, etc.). The \texttt{nanochat\_ref} field picks which fork branch or tag to run. No extra config framework needed---nanochat already parses its own arguments.

%% ============================================================
\subsection{Pipeline Architecture}
%% ============================================================

The pipeline reads a YAML config, sends training and evaluation jobs to Modal GPUs, and writes results to a Modal Volume and W\&B. All experiments are controlled via config---no pipeline code changes needed between runs.

\subsubsection{Code Structure}

\begin{forest}
  for tree={
    font=\ttfamily\small,
    grow'=0,
    child anchor=west,
    parent anchor=south,
    anchor=west,
    calign=first,
    edge path={
      \noexpand\path [draw, -{|}] (!u.south west) +(7.5pt,0) |- node[fill,inner sep=1.25pt] {} (.child anchor)\forestoption{edge label};
    },
    before typesetting nodes={
      if n=1 {insert before={[,phantom]}} {}
    },
    fit=band,
    before computing xy={l=15pt},
    s sep=4pt,
  }
  [a3/shared/scripts/
    [main.py \normalfont{\textit{--- parses YAML, calls} \texttt{run.remote()}}]
    [shared/
      [infra.py \normalfont{\textit{--- app, image, volume, secrets, setup()}}]
      [helpers.py \normalfont{\textit{--- pure Python utilities}}]
    ]
    [runners/
      [run.py \normalfont{\textit{--- orchestrator: spawns stages then evals}}]
      [train.py \normalfont{\textit{--- one training stage per GPU}}]
      [evaluate.py \normalfont{\textit{--- one checkpoint eval per GPU}}]
    ]
  ]
\end{forest}

\subsubsection{Execution Flow}

\begin{figure}[H]
\centering
\includegraphics[width=\linewidth]{pipeline_diagram.png}
\caption{Pipeline architecture. A YAML config is parsed locally, then dispatched to Modal cloud for training and evaluation. Training stages write checkpoints to a persistent Volume and metrics to W\&B; evaluations read checkpoints from the Volume and write results back.}
\label{fig:pipeline}
\end{figure}

The workflow is two steps: (1)~write a YAML config specifying model depth, sequence length, GPU type, and which evals to run, then (2)~run a single command:

\begin{lstlisting}
modal run --detach a3/shared/scripts/main.py \
    --config a3/<part>/configs/<your-config>.yaml
\end{lstlisting}

Modal spins up GPU containers, trains the model, runs evals, and writes results to a Modal Volume (checkpoints, CSVs, JSONs) and W\&B (training curves). The \texttt{--detach} flag keeps the pipeline running on Modal independent of the local terminal.

\subsubsection{Nanochat Fork}

We maintain a fork of Karpathy's nanochat at \url{https://github.com/seonghyunban/nanochat} so that teammates can modify the architecture without touching the original code. The pipeline integrates the fork as follows:

\begin{enumerate}[nosep]
  \item The Modal container image clones the fork once at build time and caches it.
  \item Before every run, the container does \texttt{git fetch} then checks out whichever branch or tag the YAML config's \texttt{nanochat\_ref} field specifies.
  \item No container rebuild is needed when new code is pushed---the next run picks it up automatically.
\end{enumerate}

\paragraph{Branch strategy}

\begin{itemize}[nosep]
  \item \texttt{master} --- original nanochat, frozen. All baseline runs use this.
  \item \texttt{p2} --- P2's architecture modifications (SwiGLU, LayerScale), branched from \texttt{master}.
  \item \texttt{p3-log-every} --- P3's fork patch adding a \texttt{--log-every} flag for finer W\&B logging.
  \item \texttt{baseline-v0} --- immutable tag on \texttt{master}. Safety net: even if someone pushes to \texttt{master}, this tag never moves.
\end{itemize}
